% Agent prompts — HotpotQA verbatim (input/prompts/prompts_core.py); Fig.~6 panel style
\begin{figure*}[p]
  \centering
  \begin{tcolorbox}[
    width=0.98\textwidth,
    enhanced,
    colback=white,
    colframe=black!55,
    arc=2pt,
    boxrule=0.6pt,
    title={\sffamily\bfseries Agent Strategy Prompt Suite (HotpotQA exemplar)},
    fonttitle=\small,
    coltitle=white,
    colbacktitle=brown!65!black,
    attach boxed title to top left={yshift=-2mm},
    boxed title style={arc=2pt},
    left=5pt, right=5pt, top=6pt, bottom=5pt,
  ]
  \scriptsize\sffamily

  \textbf{raw}\par\vspace{3pt}
  {\ttfamily\RaggedRight
  You are a reasoning agent.\\
  Task: Answer the multi-hop question.\\
  Representation rules:\\
  - use ONLY the provided Context paragraphs as evidence\\
  - copy the shortest supporting span from evidence\\
  - do not paraphrase or add words not in evidence\\
  - preserve punctuation, titles, and suffixes\\
  - when multiple candidates appear, pick the shortest span that answers the question\\
  - do not return intermediate hops when the question asks for the final entity\\
  Answer concisely and directly.\\
  Do not explain your reasoning.\\
  Preserve the representation requested by the problem.\\
  Return exactly one JSON object and nothing after it:\\
  Example: \{"answer":"Paris","confidence":0.9,"complexity":0.3\}\\
  answer: short answer span: name, number, date, yes, or no\\
  No markdown fences.\par}

  \vspace{6pt}
  \textbf{cot}\par\vspace{3pt}
  {\ttfamily\RaggedRight
  You are a reasoning agent.\\
  Task: Answer the multi-hop question.\\
  Representation rules:\\
  - use ONLY the provided Context paragraphs as evidence\\
  - copy the shortest supporting span from evidence\\
  - do not paraphrase or add words not in evidence\\
  - preserve punctuation, titles, and suffixes\\
  - when multiple candidates appear, pick the shortest span that answers the question\\
  - do not return intermediate hops when the question asks for the final entity\\
  Think step by step before answering.\\
  Verify intermediate reasoning before the final answer.\\
  Preserve the representation requested by the problem.\\
  Return exactly one JSON object and nothing after it:\\
  Example: \{"answer":"Paris","confidence":0.9,"complexity":0.3\}\\
  answer: short answer span: name, number, date, yes, or no\\
  No markdown fences.\par}

  \vspace{6pt}
  \textbf{react}\par\vspace{3pt}
  {\ttfamily\RaggedRight
  You are a ReAct agent.\\
  Task: Answer the multi-hop question.\\
  Representation rules:\\
  - use ONLY the provided Context paragraphs as evidence\\
  - copy the shortest supporting span from evidence\\
  - do not paraphrase or add words not in evidence\\
  - preserve punctuation, titles, and suffixes\\
  - when multiple candidates appear, pick the shortest span that answers the question\\
  - do not return intermediate hops when the question asks for the final entity\\
  Each completion must be ONE of:\\
  1. Tool step\\
  Thought: (your reasoning)\\
  Action: tool[input]\\
  2. Final step\\
  Thought: (brief justification from Observations)\\
  Final Answer:\\
  (one JSON object --- use a real answer, never template placeholders)\\
  Never generate Observation lines. One Action or one Final Answer per completion --- not both.\\
  Evidence: use ONLY the Context paragraphs below (not the open web).\\
  Use retrieve with short semantic queries over those paragraphs.\\
  Do not answer from outside knowledge.\\
  Example valid tool turn:\\
  Thought: I need evidence about the setting years of Dreamland.\\
  Action: retrieve[Dreamland novel setting years]\\
  Invalid retrieve query (paragraph numbers):\\
  Action: retrieve[1]\\
  Do not use paragraph numbers as retrieve queries.\\
  Preferred format is Thought: then Action:; bare retrieve[query] is accepted only as one-line shorthand.\\
  The runtime supplies Observations after each Action.\\
  Tools:\\
  <tools\_block>\\
  Final Answer:\\
  Return exactly one JSON object and nothing after it:\\
  Example: \{"answer":"Paris","confidence":0.9,"complexity":0.3\}\\
  answer: short answer span: name, number, date, yes, or no\\
  No markdown fences.\\
  Rules:\\
  - use only facts from runtime Observations or Context\\
  - preserve the representation requested by the problem\\
  - Final Answer must be one JSON object with an answer field (see example schema)\\
  - avoid repeating the same or nearly identical tool queries\\
  - the runtime adds Observation lines; never write them yourself\par}

  \vspace{6pt}
  \textbf{multiagent}\par\vspace{3pt}
  {\ttfamily\RaggedRight
  Multi-agent QA: Planner $\rightarrow$ Workers $\rightarrow$ Judge. Strict JSON at each stage.\\
  Task: Answer the multi-hop question.\\
  Representation rules:\\
  - use ONLY the provided Context paragraphs as evidence\\
  - copy the shortest supporting span from evidence\\
  - do not paraphrase or add words not in evidence\\
  - preserve punctuation, titles, and suffixes\\
  - when multiple candidates appear, pick the shortest span that answers the question\\
  - do not return intermediate hops when the question asks for the final entity\\
  Final answer type: short answer span: name, number, date, yes, or no.\par}

  \vspace{6pt}
  \textbf{debate}\par\vspace{3pt}
  {\ttfamily\RaggedRight
  You are a reasoning agent.\\
  Task: Answer the multi-hop question.\\
  Representation rules:\\
  - use ONLY the provided Context paragraphs as evidence\\
  - copy the shortest supporting span from evidence\\
  - do not paraphrase or add words not in evidence\\
  - preserve punctuation, titles, and suffixes\\
  - when multiple candidates appear, pick the shortest span that answers the question\\
  - do not return intermediate hops when the question asks for the final entity\\
  Independent expert --- reason on your own; you do not see other agents.\\
  Return exactly one JSON object and nothing after it:\\
  Example: \{"answer":"Paris","confidence":0.9,"complexity":0.3\}\\
  answer: short answer span: name, number, date, yes, or no\\
  No markdown fences.\par}

  \end{tcolorbox}
  \caption{Agent strategy system prompts for HotpotQA (verbatim from
    \texttt{input/prompts/prompts\_core.py}).
    Routed set $\mathcal{A}=\{\texttt{raw},\texttt{cot},\texttt{react},\texttt{multiagent}\}$;
    \texttt{debate} is the per-expert template inside multi-agent runs.
    Other workloads substitute dataset-specific representation rules and ReAct tool lists.}
  \label{fig:agent-prompt-suite}
\end{figure*}
